% =====================================================================
\section{Цифровой вход: кнопки и прерывания}
\label{les:gpio-in}
% =====================================================================
\lessonintro
  {научиться читать кнопку, понять роль подтягивающего резистора, победить дребезг, познакомиться с прерываниями}
  {урок~\ref{les:signals} (логические уровни), урок~\ref{les:led} (\code{millis()}), урок~\ref{les:pwm} (яркость)}
  {кнопка переключает яркость светодиода: выкл → 30\,\% → 100\,\%}
  {тактовая кнопка, светодиод с резистором \SI{220}{\ohm}}

До сих пор выводы платы только \emph{выдавали} сигналы. Теперь научимся их \emph{читать}. Самое
простое устройство ввода --- кнопка: нажата или нет, 1 или 0. Идеальный цифровой сигнал!

\subsection{Как устроена кнопка}

Внутри тактовой кнопки две металлические пластинки. Когда кнопку нажимают, пластинки касаются, и цепь
замыкается. У кнопки четыре ножки, но соединены они попарно: ножки на противоположных сторонах
всегда связаны, а нажатие соединяет одну пару с другой.

\begin{figure}[H]
\centering
\begin{tikzpicture}[x=1cm,y=1cm]
  \fill[black!75,rounded corners=1mm] (0,0) rectangle (2,2);
  \fill[black!40] (1,1) circle (0.55);
  \foreach \p/\l in {{(-0.3,0.3)}/1,{(2.3,0.3)}/2,{(-0.3,1.7)}/3,{(2.3,1.7)}/4}{\fill[black!55] \p ++(-0.15,-0.15) rectangle ++(0.3,0.3); }
  \node[font=\sffamily\scriptsize] at (-0.3,-0.2) {A}; \node[font=\sffamily\scriptsize] at (2.3,-0.2) {B};
  \node[font=\sffamily\scriptsize] at (-0.3,2.2) {A}; \node[font=\sffamily\scriptsize] at (2.3,2.2) {B};
  \node[font=\sffamily\small,align=left,anchor=west] at (3,1) {ножки A соединены всегда,\\ ножки B --- тоже;\\ нажатие соединяет A с B};
  \begin{scope}[xshift=9cm,yshift=1cm]
    \draw (0,0) node[left,font=\sffamily\small]{A} to[push button,o-o] (2,0) node[right,font=\sffamily\small]{B};
  \end{scope}
\end{tikzpicture}
\caption{Тактовая кнопка (вид снизу) и её обозначение на схеме}
\end{figure}

\subsection{Проблема «висящего» входа}

Кажется, всё просто: подключить кнопку между выводом и землёй. Нажали --- на выводе 0. А если
отпустили? Тогда вывод вообще ни к чему не подключён. Вход ESP32-S3 очень чувствительный: «висящий»
провод работает как антенна и ловит наводки от рук, телефона, проводки. Программа будет читать то 0,
то 1 --- случайно.

Решение --- \textbf{подтягивающий резистор}. Он слабо «тянет» вывод к 3,3~В, и в покое там
уверенная 1. Нажатая кнопка соединяет вывод с землёй гораздо сильнее, чем резистор, --- и получается 0.

\begin{figure}[H]
\centering
\begin{circuitikz}
  \draw (0,3) node[vcc]{3,3 В} to[R,l_=$R$ \SI{45}{k\ohm}] (0,1) coordinate(a);
  \draw (a) to[short,-o] (1.5,1) node[right]{\pin{5}};
  \draw (a) to[push button, l_=S1] (0,-1) node[ground]{};
  \node[align=center,font=\sffamily\small] at (0.6,-2.3) {Подтяжка к питанию (pull-up):\\ отпущена = 1, нажата = 0};
  \begin{scope}[xshift=7cm]
  \draw (0,3) node[vcc]{3,3 В} to[push button, l_=S1] (0,1) coordinate(b);
  \draw (b) to[short,-o] (1.5,1) node[right]{\pin{5}};
  \draw (b) to[R,l_=$R$ \SI{45}{k\ohm}] (0,-1) node[ground]{};
  \node[align=center,font=\sffamily\small] at (0.6,-2.3) {Подтяжка к земле (pull-down):\\ отпущена = 0, нажата = 1};
  \end{scope}
\end{circuitikz}
\caption{Два способа подтяжки. У ESP32-S3 оба резистора уже есть внутри чипа, их включают программой.
Мы будем использовать pull-up: кнопка замыкает вывод на землю}
\end{figure}

\subsection{Схема подключения}

Добавим к светодиоду на \pin{4} кнопку на \pin{5}. Внешний резистор для кнопки не нужен: включим
встроенную подтяжку командой \code{pinMode(5, INPUT\_PULLUP)}.

\begin{twoview}{Кнопка на GPIO5 и светодиод на GPIO4}{fig:button}
\begin{tikzpicture}[bb]
  \bbBoard{24}
  \bbDevKit{29}{25}
  \bbR{10}{13}{14}{13}{rr}{rr}{rn}
  \bbLED{red}{14}{11}{-1.8}
  \draw[wire=wblack] (15,14) -- (15,16);
  \draw[wire=wyellow] (dk-4) -- (10,22) -- (10,14);
  \bbButton{18}
  \draw[wire=wgreen] (dk-5) -- (18,21) -- (18,14);
  \draw[wire=wblack] (20,14) -- (20,16);
  \draw[wire=wblack] (dk-GND) -- (26.5,4) -- (26.5,16) -- (24,16);
\end{tikzpicture}
\nextview
\begin{circuitikz}
  \draw (0,0) node[left]{\pin{4}} to[short,o-] (0.5,0)
        to[R, l=\SI{220}{\ohm}] (3,0) to[led] (5.2,0) -- (5.8,0) -- (5.8,-0.5) node[ground]{};
  \draw[dashed,gray] (-2.7,-1.4) rectangle (-0.3,-4.6);
  \node[font=\sffamily\scriptsize,text=gray] at (-1.5,-4.35) {внутри чипа};
  \draw (-1.2,-2.2) node[vcc]{3,3 В} to[R,l_=\SI{45}{k\ohm}] (-1.2,-3.6) -- (0,-3.6);
  \draw (0,-3.6) to[short,o-] (0.5,-3.6) to[push button,l=S1] (3,-3.6) -- (5.8,-3.6) -- (5.8,-4.1) node[ground]{};
  \node[font=\sffamily\bfseries\small,below] at (0,-3.7) {GPIO5};
\end{circuitikz}
\end{twoview}

На макетной плате кнопка стоит поперёк канавки. Зелёный провод от \pin{5} подключён к одной её стороне,
чёрный провод от другой стороны ведёт на шину земли.

\subsection{Читаем кнопку}

\codecap{Состояние кнопки в Serial Monitor}
\codeinput{l10_button_read}

Попробуй закомментировать \code{INPUT\_PULLUP} и написать просто \code{INPUT}. Поднеси руку к проводу
--- увидишь, как «висящий» вход сходит с ума.

\subsection{Дребезг контактов}

Если просто написать «когда кнопка нажата --- переключи светодиод», он будет переключаться много раз
за одно нажатие. Виноват \textbf{дребезг}: металлические пластинки при ударе несколько раз
«подпрыгивают», как мячик. Для человека это одно нажатие, а микроконтроллер, проверяющий вход
миллионы раз в секунду, видит пачку нажатий.

\begin{figure}[H]
\centering
\begin{tikzpicture}[x=1cm,y=1cm]
  \draw[->] (0,0) -- (12.3,0) node[right]{$t$};
  \draw[->] (0,0) -- (0,2) node[above]{$U$};
  \draw (0.05,1.4) -- (-0.05,1.4) node[left,font=\scriptsize]{1};
  \draw (0.05,0) -- (-0.05,0) node[left,font=\scriptsize]{0};
  \draw[very thick,accent] (0,1.4) -- (2,1.4) -- (2,0) -- (2.2,0) -- (2.2,1.2) -- (2.35,1.2) -- (2.35,0)
    -- (2.6,0) -- (2.6,0.9) -- (2.7,0.9) -- (2.7,0) -- (2.95,0) -- (2.95,0.5) -- (3.0,0.5) -- (3.0,0) -- (8,0)
    -- (8,1.4) -- (8.12,1.4) -- (8.12,0.3) -- (8.3,0.3) -- (8.3,1.4) -- (8.45,1.4) -- (8.45,0.8)-- (8.55,0.8) -- (8.55,1.4) -- (12,1.4);
  \draw[decorate,decoration={brace,amplitude=4pt}] (2,1.6) -- (3.0,1.6) node[midway,above=4pt,font=\scriptsize,align=center]{дребезг\\ 1--20 мс};
  \draw[decorate,decoration={brace,amplitude=4pt}] (8,1.6) -- (8.55,1.6) node[midway,above=4pt,font=\scriptsize]{дребезг};
  \node[font=\scriptsize] at (5.5,0.35) {кнопка нажата};
  \begin{scope}[yshift=-2.4cm]
  \draw[->] (0,0) -- (12.3,0) node[right]{$t$};
  \draw[->] (0,0) -- (0,2);
  \node[font=\scriptsize,anchor=south west] at (0.2,1.5) {после фильтра};
  \draw[very thick,okgreen] (0,1.4) -- (4.0,1.4) -- (4.0,0) -- (9.55,0) -- (9.55,1.4) -- (12,1.4);
  \draw[dashed,gray] (3.0,0) -- (3.0,4.0);
  \draw[dashed,gray] (8.55,0) -- (8.55,4.0);
  \draw[<->,esp] (3.0,0.7) -- node[above,font=\scriptsize]{50 мс} (4.0,0.7);
  \draw[<->,esp] (8.55,0.7) -- node[above,font=\scriptsize]{50 мс} (9.55,0.7);
  \node[font=\scriptsize,text=okgreen] at (4.0,-0.35) {событие «нажатие»};
  \end{scope}
\end{tikzpicture}
\caption{Дребезг контактов (сверху) и сигнал после программного фильтра (снизу)}
\end{figure}

Программный фильтр верит новому уровню, только если тот продержался дольше 50~мс. Оформим его в функцию,
которая возвращает \code{true} ровно один раз на каждое нажатие:

\codecap{Функция buttonPressed(): нажатие без дребезга}
\codeinput{l10_debounce}

Слово \code{static} перед переменной внутри функции означает: «помни значение между вызовами».

\subsection{Практика: кнопка-переключатель яркости}

Каждое нажатие переключает светодиод на следующую ступень: выключен → 30\,\% → 100\,\% → выключен.
Так работают многие настольные лампы.

\begin{figure}[H]
\centering
\begin{tikzpicture}[node distance=15mm,st/.style={blk,circle,minimum size=17mm,fill=blkmem}]
  \node[st,fill=gray!20] (a) {выкл\\ 0\,\%};
  \node[st,fill=esp!20,right=of a] (b) {30\,\%};
  \node[st,fill=esp!55,right=of b] (c) {100\,\%};
  \draw[arr] (a) -- node[above,font=\sffamily\scriptsize]{нажатие} (b);
  \draw[arr] (b) -- node[above,font=\sffamily\scriptsize]{нажатие} (c);
  \draw[arr] (c) to[bend left=35] node[below,font=\sffamily\scriptsize]{нажатие} (a);
\end{tikzpicture}
\caption{Диаграмма состояний: кружки --- состояния, стрелки --- что их меняет}
\end{figure}

\codecap{Кнопка переключает яркость}
\codeinput{l10_toggle}

\subsection{Прерывания}

Пока \code{loop()} крутится быстро, опрос кнопки работает отлично. Но если программа занята
чем-то долгим, короткое нажатие можно пропустить. Выход --- \textbf{прерывание}: попросить процессор
\emph{самого} вызвать нашу функцию в тот момент, когда уровень на выводе изменится. Как будильник:
не нужно всё время смотреть на часы, он сам позовёт.

\begin{figure}[H]
\centering
\begin{tikzpicture}[x=1cm,y=1cm]
  \draw[thick,->] (0,0) -- (12,0) node[right,font=\small]{основная программа};
  \foreach \x in {0.3,1.3,...,11.3} \fill[accent!40] (\x,-0.15) rectangle (\x+0.8,0.15);
  \draw[very thick,esp,->] (4.7,1.3) node[above,font=\small]{нажата кнопка} -- (4.7,0.2);
  \fill[esp!40] (4.8,-1.2) rectangle (6,-0.8) node[midway,font=\scriptsize]{ISR};
  \draw[->,esp] (4.7,-0.15) -- (4.8,-1.0);
  \draw[->,esp] (6,-1.0) -- (6.1,-0.15);
  \node[font=\scriptsize,align=center] at (8.5,-1.1) {обработчик прерывания:\\ коротко! отметить и выйти};
\end{tikzpicture}
\caption{Прерывание ненадолго приостанавливает основную программу}
\end{figure}

\codecap{Подсчёт нажатий через прерывание}
\codeinput{l10_interrupt}

\begin{caution}
Внутри обработчика прерывания (ISR) нельзя вызывать \code{delay()}, \code{Serial.print()} и другие медленные
функции. Правило: изменить переменную --- и выйти. Всю остальную работу делает основная программа.
\end{caution}

\begin{tip}
Кнопку не стоит подключать к «загрузочным» выводам \pin{0}, \pin{3}, \pin{45}, \pin{46} (урок~\ref{les:pins}):
если при включении она окажется нажатой, плата может не запустить программу.
\end{tip}

\begin{summary}
\begin{itemize}[leftmargin=*]
  \item Вход без подтяжки «висит». Используем \code{INPUT\_PULLUP} и кнопку на землю: нажата = \code{LOW}.
  \item Ловим \emph{момент} нажатия и отфильтровываем дребезг (20--50~мс).
  \item Прерывание срабатывает мгновенно, но обработчик должен быть очень коротким.
\end{itemize}
\end{summary}

\begin{quiz}
\begin{enumerate}
  \item Что прочитает программа с вывода с \code{INPUT\_PULLUP}, если кнопка не нажата?
  \item Почему одно нажатие может посчитаться несколько раз?
  \item Чего нельзя делать внутри обработчика прерывания?
\end{enumerate}
\end{quiz}

\begin{tasks}
\begin{enumerate}
  \item Используй кнопку BOOT (\pin{0}) как вторую кнопку --- после того как программа уже запустилась.
  \item Различай короткое (< 0,5~с) и длинное (> 1~с) нажатие: длинное выключает светодиод сразу.
  \item Замени кнопку сенсорным входом: \code{touchRead(2)} растёт, если коснуться провода на \pin{2}.
  \item Сделай «реакциометр»: светодиод загорается в случайный момент, программа измеряет, как быстро ты нажмёшь кнопку.
\end{enumerate}
\end{tasks}
